\section{The finite nonlinear comparison}\label{sec:dn-comparison}

We now evolve the actual entry data supplied by
\Cref{prop:packet-preparation}.  All comparisons in this Section use the
same numerical optical and angular labels.  A bar denotes the exact
double-null template, a prime denotes the vacuum solution to be
constructed, and an undecorated field belongs to the fixed background.
In particular, the barred metric is generally not vacuum.  Its equations
have the controlled defects provided by \Cref{prop:packet-preparation}.
For each decoration the metric is
\[
 -2q\mathcal A\,\dd u\,\dd v+
   \gamma_{AB}(\dd\theta^A-b^A\dd u)(\dd\theta^B-b^B\dd u).
\]
The function $q$ and the constants $q_0,\kappa$ are common to all
three descriptions; it is the positive relative factor $\mathcal A$
that changes.

Fix a finite desired output order $s$ and a finite initial-neighborhood
order $m$.  Enlarging $s$ if necessary, we assume $s\ge30$.  The domain
from \Cref{prop:packet-preparation} is
\begin{equation}\label{dn:domain}
\begin{split}
 \Omega_T=\{(t,x,\theta):\;&0\le t\le L_T=(T+u_*)/2+c_0,\\
 &u=t-x\le u_2-c_0t/T,\quad
 v=t+x\le T+c_2-c_0t/T,\quad\theta\in\Sph\}.
\end{split}
\end{equation}
The fixed constants and their small enlargements are chosen as in
\Cref{lem:bg-branch-tube}, so the observation tube has an interior
margin and every spatial interval $I_T(t)$ has length at least a fixed
positive number.  Its length is at most $C(T+1)$.  We use these
enlargements when taking traces and restarting at an endpoint.  Put
\[
 q=q_0e^{-\kappa(u+v)},\qquad \sigma=\sqrt q,
 \qquad \alpha=\kappa/4,\qquad \amp=e^{-\alpha T}.
\]
Neither the constants below nor the orders chosen below require a
further reduction of the fixed neighborhood $\Ueps$.

\begin{theorem}[Finite double-null experiment]\label{thm:dn-comparison}
For every fixed background and branch tube in
\Cref{lem:bg-branch-tube}, every finite $m,s$, and every sufficiently
small fixed $c>0$, the parameters in \Cref{prop:packet-preparation} can
be chosen so that its exact entry data have a smooth vacuum development
on \Cref{dn:domain} in the double-null gauge initialized there.
The relative lapse and the sphere metric remain positive and uniformly
nondegenerate.  For
\[
 \mathcal U=(\gamma,\log\mathcal A,b,\chi_0,\chi_1,
                       \ell_0,\ell_1,\xi_0,\xi_1)
\]
one has, in a fixed sphere atlas and ordinary coordinate derivatives,
\begin{equation}\label{dn:main-error}
 \max_{|I|\le s}\|\partial^I(\mathcal U'-\bar{\mathcal U})
                         \|_{L^\infty(\Omega_T)}
       \le e^{-(1+c)\alpha T}
\end{equation}
for all sufficiently large $T$ in the selected experiment sequence.
The same conclusion holds for
$q^{-1}\partial_vb'-q^{-1}\partial_v\bar b$ at any preassigned finite
order, after increasing the finite order budget.  The individual
ordinary jets of the primed relative fields through every order used
in the experiment are $e^{o(T)}$, uniformly for this moving family.
The exact solution and the template equal the background where
$v<(1-\nu)T$.  The exact solution agrees on an entry collar, by local
geometric uniqueness, with the actual exterior solution of
\Cref{prop:packet-preparation}.
\end{theorem}

Three estimates connect the prepared entry data to this conclusion.
Under a low-order bootstrap, the individual exact fields first satisfy
high-order bounds with arbitrarily small exponential loss. We then prove
a quadratic $L^2$ bound for the actual-minus-template difference, keeping
the lapse factors that compensate the angular elliptic estimates.
Interpolation between these two bounds yields the required finite-jet
error and improves the bootstrap. The estimates are proved on every
existing partial slab, so a local restart completes the existence
argument. The quadratic bound is \Cref{prop:dn-low-error}; the next
subsection begins with the equations needed for both estimates.

\subsection{Equations, bootstrap, and finite-order data}

For any of the three metrics define, with $i\ne j$,
\begin{equation}\label{dn:definitions}
\begin{gathered}
 a=q\mathcal A,\quad B_0=b,\quad B_1=0,\quad
 L_i=\partial_{y_i}+B_i,\quad
 \mathfrak L_i=\partial_{y_i}+\Lie_{B_i},\quad (y_0,y_1)=(u,v),\\
 \chi_i=\tfrac12\mathfrak L_i\gamma,\quad
 \tau_i=\tr_\gamma\chi_i,\quad \ell_i=L_i\log a,\quad
 \xi_0+\xi_1=\dd_\theta\log\mathcal A,\quad
 \xi_0-\xi_1=-a^{-1}\gamma\partial_vb,\\
 \beta_i=\divg_\gamma\chi_i-\dd_\theta\tau_i,\qquad
 c_i=\curl_\gamma\xi_j,\qquad
 Q_i=(K_\gamma,c_i),\qquad
 m_i=K_\gamma-\divg_\gamma\xi_j.
\end{gathered}
\end{equation}
Here $K_\gamma$ is Gaussian curvature, and
$\operatorname{Sym}T=(T+T^{\mathsf t})/2$.  Products of shape tensors
contract their adjacent indices using $\gamma^{-1}$.
The vacuum double-null equations of
\cite[\CompanionRef{foundation}{in:vacuum}, \CompanionRef{foundation}{in:angular-system}]{KerrCTwo} read
\begin{equation}\label{dn:D2}
\begin{aligned}
 L_i\tau_i&=\ell_i\tau_i-|\chi_i|_\gamma^2,\\
 \mathfrak L_j\chi_i
 &=a\operatorname{Sym}\nabla\xi_j-\tfrac a2K_\gamma\gamma
       +a\xi_j^{\otimes2}
       -\tfrac12(\tau_i\chi_j+\tau_j\chi_i)
       +\chi_i\chi_j+\chi_j\chi_i,\\
 \mathfrak L_i\xi_j&=\beta_i+\tau_i\xi_i,\\
 L_j\ell_i&=aK_\gamma-\chi_i\cdot\chi_j+\tau_i\tau_j
                   +a(2\xi_i\cdot\xi_j-|\xi_i|^2).
\end{aligned}
\end{equation}
The exact definition identities include
\begin{equation}\label{dn:recovery-equations}
 \partial_v\gamma=2\chi_1,\qquad
 \partial_v\log\mathcal A=\ell_1+\kappa,\qquad
 \partial_vb=q\mathcal A\gamma^{-1}(\xi_1-\xi_0).
\end{equation}
Combining Raychaudhuri with
$L_i\log\mathcal A=\ell_i+\kappa$ gives the particularly useful identity
\begin{equation}\label{dn:damped-trace}
 (L_i+\kappa)w_i=-\mathcal A^{-1}|\chi_i|^2,
                  \qquad w_i=\tau_i/\mathcal A.
\end{equation}
Thus the constant damping survives division by the relative lapse.

For clarity, the angular system is also recorded.  The symbol $*$
means a finite sum of smooth contractions with $\gamma^{\pm1}$,
with precisely the differentiated factors displayed; it conceals no
additional derivative.  With $J\dd f=(-\partial_2f,\partial_1f)$ in
oriented orthonormal components, one has
\begin{equation}\label{dn:D3}
\begin{aligned}
 \mathfrak L_j\beta_i
 &=\tfrac a2(\dd K_\gamma+J\dd c_i)
       +a(K_\gamma*\xi+\xi*\nabla\xi+\xi^3)
       +\chi_i*\nabla\chi_j+\chi_j*\nabla\chi_i,\\
 L_iQ_i&=(\divg,\curl)\beta_i+\mathcal R_i,
       &L_im_i&=\mathcal T_i,\\
 \mathcal R_i,\mathcal T_i
 &:\quad\chi_i*(K_\gamma,\nabla\xi)+\nabla\chi_i*\xi.
\end{aligned}
\end{equation}
The factor list can be checked without differentiating a spacetime
curvature component: apply divergence minus trace differential to the
shape cross equation, and use
\begin{align*}
 L_iK_\gamma&=\divg\beta_i-\tau_iK_\gamma,\\
 L_i\divg\zeta
 &=\divg(\mathfrak L_i\zeta)
      -2\chi_i^{AB}\nabla_A\zeta_B
      -(2\divg\chi_i-\dd\tau_i)\cdot\zeta,\\
 L_i\curl\xi_j
 &=\curl\beta_i+\curl(\tau_i\xi_i)-\tau_i\curl\xi_j.
\end{align*}
In particular, the divergence of $\beta_i$ cancels from $L_im_i$.
More explicitly,
\[
 L_im_i=-\tau_iK_\gamma-\divg(\tau_i\xi_i)
       +2\chi_i^{AB}\nabla_A(\xi_j)_B
       +(2\divg\chi_i-\dd\tau_i)\cdot\xi_j,
\]
which verifies all the displayed factor types of its source.
These are the identities used to derive the angular system in
\cite[\CompanionRef{foundation}{in:angular-system}]{KerrCTwo}.

We work first on an arbitrary existing partial slab.  Bootstrap the
ordinary derivatives of $\mathcal U'-\mathcal U$ through order ten by
\begin{equation}\label{dn:bootstrap}
 \|\mathcal U'-\mathcal U\|_{C^{10}}\le e_T,
                   \qquad e_T=e^{-\alpha T/2},
\end{equation}
and keep the relative lapse and sphere metric in a fixed open
nondegeneracy range about the background bounds.  One may enlarge the
fixed bootstrap order when needed for a local existence convention.
\Cref{prop:bg-package} and \Cref{dn:bootstrap} imply uniform angular
$H^6$ bounds for $\gamma'^{\pm1},\mathcal A'^{\pm1}$ and $H^5$ bounds
for the torsions.  All required low angular norms of
$\chi_i',\ell_i'+\kappa,B_i'$ and the low pure longitudinal
coefficient norms are majorized by
\begin{equation}\label{dn:majorants}
 h_{i,T}=C_0(q+g_i(y_i)+e_T).
\end{equation}
The functions $g_i\ge0$ are fixed background envelopes with summable
unit-cell suprema, enlarged to include the pure-jet controls of
\cite[\CompanionRef{foundation}{in:pure-jets}]{KerrCTwo}.  The barred fields obey the same
bounds after the frequency growth exponent $\zeta$ is chosen
sufficiently small at the finitely many base orders.  In the rest
of the Section $h_i$ denotes $h_{i,T}$.

Given any finite order $N$ and positive loss $\eta$, the preparation
provides the initial discrepancy and all required defects of
\Cref{dn:D2,dn:D3,dn:damped-trace} bounded by
\begin{equation}\label{dn:input-defect}
 F_T=\amp^2e^{\eta T}
\end{equation}
through that order.  The geometric definitions and
\Cref{dn:recovery-equations} hold exactly for the template.  Its
finite jets differ from the background by
$\amp\lambda^{P_N}e^{o(T)}$, including the relative shift derivative.
The positive frequency exponent $P_N$ is fixed by the base orders,
independently of subsequently requested inverse-frequency accuracy.
The initial high norms needed below are consequently bounded by
$e^{o(T)}$ uniformly in $T$ once $\zeta P_N<\alpha/2$.
Only actual entry data, rather than the template, initialize the
vacuum Cauchy problem.

\subsection{Uniform high estimates for the moving experiments}

\begin{lemma}[Uniform finite-slab high bounds]\label{lem:dn-high-uniform}
Assume \Cref{dn:bootstrap} on an existing initial partial slab
$\Omega_T\cap\{0\le t\le t_*\}$ of \Cref{dn:domain}, with the same
outer margins.  For every fixed finite order $n$ there is a finite
initial order $r(n)$ with the following property.  Uniform
$e^{o(T)}$ initial bounds through order $r(n)$ imply uniform
$e^{o(T)}$ bounds through order $n$ for all ordinary jets of
\[
 \gamma'^{\pm1},\quad\log\mathcal A',\quad b',\quad
 \chi_i',\quad\ell_i',\quad\xi_i',\quad
                    q^{-1}\partial_vb'.
\]
The constants depend only on the fixed orders, a prescribed
exponential-rate tolerance, the bounded low coefficient set, the
uniform envelope bounds described below, and finitely many initial
norms.  They do not depend on an uncontrolled high norm of the
unknown solution or on the subslab endpoint.
\end{lemma}

\begin{proof}
We give the finite-subslab argument, including the dependence needed
here.  We reproduce the finite-slab argument of
\cite[\CompanionRef{foundation}{in:higher}, proof]{KerrCTwo}, with uniformly
controlled envelopes.  Merely applying the fixed-datum conclusion of
that Proposition would not establish the present assertion.

Choose smooth cutoffs $0\le\psi_{i,T}\le1$ equal to one on the coordinate projection
of $\Omega_T$, supported in $[-C(T+1),C(T+1)]$.  Extend only the
scalar envelope, not the spacetime fields, and set
\begin{equation}\label{dn:uniform-envelopes}
 G_{i,T}(y)=g_i(y)+e_T\psi_{i,T}(y),
                       \qquad h_i=C_0(q+G_{i,T}).
\end{equation}
The added $L^1$ mass and sum of unit-cell suprema are
$O((T+1)e_T)$, uniformly bounded and tending to zero.  Their tails
are uniformly tight: if the added support reaches $|y|>R$, then
$T+1\ge R/C$, and hence
\[
 \sup_T\int_{|y|>R}e_T\psi_{i,T}(y)\,\dd y
 \le C\sup_{T+1\ge R/C}(T+1)e^{-\alpha T/2}\longrightarrow0.
\]
The same argument controls the tail supremum.  Together with the
fixed summable $g_i$, this supplies all tail thresholds uniformly.
The Lipschitz potential
\begin{equation}\label{dn:bounded-potential}
 \Phi_T(u,v)=C_0\sum_{i=0}^1\int_{-\infty}^{y_i}G_{i,T}(r)\,\dd r
       +\frac{C_0q_0}{\kappa}(1-e^{-\kappa(u+v)})
\end{equation}
satisfies $0\le\Phi_T\le C$ and
$L_i'\Phi_T=\partial_i\Phi_T=h_i$ almost everywhere.
In particular both directional integrals of $h_i$ are uniformly
bounded on every past transport segment in the domain.

The weighted angular curvature pair will supply stable fluxes, while
trace damping controls the shape trace. We will then recover unweighted
metric and torsion fields by transport, retaining the small lapse in the
curvature energy throughout.

Suppress primes in this proof only.  Fix a smooth reference metric on
$\Sph$ and let $D$ be its Levi-Civita connection, independent of
$u,v,T$ and of the solution.  In this proof all angular Sobolev sizes
use the tensors $D^r f$; equivalently, they use component derivatives
in the fixed finite atlas of the background.  The low coefficient
bounds are used in this same reference convention, together with the
uniform positive upper and lower metric bounds in that atlas.  This is
the convention for the bounded low coefficient set below.  The symbol
$\nabla$ in the equations continues to mean the connection of
$\gamma$.

For definiteness, put
\[
 |f|_k^2=\sum_{r=0}^k\int_{\Sph}|D^r f|_\gamma^2\,\dd\mu_\gamma.
\]
The contraction includes every angular tensor slot; finite scalar
arrays have their Euclidean contraction.  These norms are uniformly
equivalent to the fixed-reference Sobolev norms at each fixed order.
In particular $D\gamma$ need not vanish, so $|\gamma|_k$ records the
coordinate coefficient jets.  At a fixed high angular order $M\ge5$,
use the actual-field sizes
\begin{equation}\label{dn:high-sizes}
\begin{gathered}
 \mathsf Y_i=|\beta_i|_M,\quad
 \mathsf Z_i=\sigma|(Q_i,m_i)|_M,\quad
 \mathsf R_i=|\chi_i|_{M+1},\quad
 \mathsf z_i=|w_i|_{M+1},\quad \mathsf d_i=|B_i|_{M+1},\\
 \mathsf d_1=0,\quad\mathsf D=\mathsf d_0,\quad
 \mathsf P_\gamma=1+|\gamma|_{M+1},\quad
 \mathsf P_i=|\xi_j|_M,\quad
 \mathsf P=\mathsf P_\gamma+\mathsf P_0+\mathsf P_1.
\end{gathered}
\end{equation}
Inverse metrics and relative lapse factors through order $M+1$ are
bounded by $C_M\mathsf P$: differentiate inversion tamely and use
$\dd\log\mathcal A=\xi_0+\xi_1$.  The undifferentiated mean of
$\log\mathcal A$ is already bounded by the low bounds for
$\mathcal A^{\pm1}$.
The angular elliptic estimates are
\begin{equation}\label{dn:high-hodge}
 \mathsf R_i\le C_M(\mathsf Y_i+\mathsf z_i+h_i\mathsf P),
 \qquad
 |\xi_j|_{M+1}\le C_M(\sigma^{-1}\mathsf Z_i+\mathsf P).
\end{equation}
To make the constant's dependence explicit, freeze the principal
coefficient in sufficiently small fixed charts, identify trace-free
tensors by smooth metric-dependent bases, and use the first-order
elliptic estimate there.  The differentiated coefficient error is
bounded by
$C_M(|f|_M+|\gamma|_{M+1}\|f\|_{C^1})$.
Interpolation absorbs the penultimate derivative of $f$ with a
constant depending on $M$ and the low Lipschitz/ellipticity bounds.
Every high metric coefficient thus remains a variable in
$\mathsf P$; it does not enter $C_M$.  In the shape estimate the
low factor multiplying that variable is $h_i$, including conversion
of $\tau_i=\mathcal A w_i$.  Equivalently, the trace-free tensor
$S_i=\chi_i-\tfrac12\tau_i\gamma$ satisfies
$\divg S_i=\beta_i+\tfrac12\dd\tau_i$, and the one-form data are
$\divg\xi_j=K_\gamma-m_i$ and $\curl\xi_j=c_i$.

Commute \Cref{dn:D3} with $D^r$, $0\le r\le M$.  Transport the
resulting angular tensor $D^r f$ by
$\mathfrak L_i(D^r f)=\partial_{y_i}D^r f+\Lie_{B_i}D^r f$ on all
of its slots.  In particular, for a scalar $f$ the commutator is
$D^rL_i f-\mathfrak L_i(D^r f)$, because its $r$ derivative slots
are now covariant tensor slots.  In the commuted principal operators,
$\nabla$ acts on every slot.  If $I$ denotes the $r$ derivative
slots, the new gradient slot and the original one-form slot are paired by
\[
 (\nabla F)_{IA}=\nabla_A F_I,\qquad
 (\divg\mathcal B)_I=\gamma^{AB}\nabla_A\mathcal B_{IB},\qquad
 (\curl\mathcal B)_I=\epsilon_\gamma^{AB}\nabla_A\mathcal B_{IB},
\]
where $\epsilon_\gamma$ is the oriented area bivector.  Use the same
contraction of the $I$ slots in both members of the pair.  Their densities are
\begin{equation}\label{dn:high-pair-densities}
\begin{aligned}
 e_{\mathsf Y_i}
   &=\frac12\sum_{r=0}^M|D^r\beta_i|_\gamma^2
       &&\text{in direction }j,\\
 e_{\mathsf Z_i}
   &=\frac a4\sum_{r=0}^M
       \bigl(|D^rQ_i|_\gamma^2+|D^rm_i|_\gamma^2\bigr)
       &&\text{in direction }i .
\end{aligned}
\end{equation}
The scalar to one-form weight ratio is exactly $a/2$, and the
$m_i$ current has no angular principal derivative.  At each order,
integration by parts on the full sphere cancels the two principal
terms except for
\[
 -\frac12\int_{\Sph}
 \big\langle D^r\beta_i,\,
       (\dd a)D^rK_\gamma+J(\dd a)D^rc_i
 \big\rangle_\gamma\,\dd\mu_\gamma .
\]
This uses $\nabla J=0$ and the same derivative-slot contraction on
both sides.  Since $\dd a=a(\xi_0+\xi_1)$, the sum of these errors is
bounded by $C_M\sigma\mathsf Y_i\mathsf Z_i$.  The sphere integrals of
the densities in
\Cref{dn:high-pair-densities} are uniformly equivalent to
$\mathsf Y_i^2,\mathsf Z_i^2$, because $a=q\mathcal A$ and
$\mathcal A^{\pm1}$ are bounded.

Here are representative commutator bounds that specify the high
coefficient dependence.  The difference between the two angular
connections is
\[
 (\nabla-D)^C{}_{AB}
 =\tfrac12\gamma^{CD}
       (D_A\gamma_{BD}+D_B\gamma_{AD}-D_D\gamma_{AB}),
 \qquad |\,\nabla-D\,|_M\le C_M\mathsf P_\gamma.
\]
Tame products, the low bounds for $\dd Q_i$, and
$|\beta_i|_{C^1}\le C h_i$ give, for $r\le M$,
\begin{equation}\label{dn:high-commutators}
\begin{aligned}
 \big|D^r\{a(\dd K_\gamma+J\dd c_i)\}
      -a\{\nabla D^rK_\gamma+J\nabla D^rc_i\}\big|_0
   &\le C_M(\sigma\mathsf Z_i+q\mathsf P),\\
 \big|D^r\{(\divg,\curl)\beta_i\}
             -(\divg,\curl)D^r\beta_i\big|_0
   &\le C_M(\mathsf Y_i+h_i\mathsf P),\\
 \big|[D^r,\mathfrak L_j]\beta_i\big|_0
   &\le C_M(h_j\mathsf Y_i+h_i\mathsf d_j).
\end{aligned}
\end{equation}
In the first line, every high derivative of $a$ or $J$ multiplies a
low derivative of $Q_i$ and costs $q\mathsf P$; the other terms cost
$q|Q_i|_M\le\sigma\mathsf Z_i$.  In the second line, a high derivative
of $\gamma^{-1}$ or $\nabla-D$ multiplies low $\dd\beta_i$ or
$\beta_i$, giving $h_i\mathsf P$; its scalar-current weight supplies
one further $\sigma$.  The fixed reference curvature contributes only
lower derivatives with these same factors.

For the last line, the identity
$\Lie_B f=B*Df+DB*f$ for a covariant tensor shows that the
commutator has low $DB_j$ against the top $\beta_i$ jet, or a
derivative of $B_j$ through order $M+1$ against a low $\beta_i$ jet.
Reference-curvature terms have the additional low factor $B_j$.
Likewise the scalar $Q_i,m_i$ commutators cost
$C_M\{\sigma\mathsf d_i+h_i\mathsf Z_i\}$ after weighting.
The order-$M+1$ trace commutator costs
$C_Mh_i(\mathsf z_i+\mathsf d_i)$, whereas the order-$M$ torsion
commutator costs $C_M(h_i\mathsf P_i+\mathsf d_i)$.
The time-independent choice of $D$ also gives
$[\partial_v,D^r]b=0$ in the shift equation.

The remaining differentiated products have the same tame placement.
A high angular connection in $\chi_i*\nabla\chi_j$ costs
$h_ih_j\mathsf P$; in $a\xi*\nabla\xi$ it costs $q\mathsf P$.
The estimate for a differentiated torsion in an $a$ term is
\[
 q\bigl(|\xi_0|_{M+1}+|\xi_1|_{M+1}\bigr)
 \le C_M\{\sigma(\mathsf Z_0+\mathsf Z_1)+q\mathsf P\}
\]
by \Cref{dn:high-hodge}.  These observations and Leibniz placement give,
apart from favorable damping, paired squared-current costs at most
$C_M$ times
\begin{equation}\label{dn:high-pair-costs}
\begin{aligned}
 \mathsf Y_i\big[&\sigma(\mathsf Z_0+\mathsf Z_1+\mathsf P)
       +h_j(\mathsf Y_i+\mathsf R_i)
       +h_i(\mathsf R_j+\mathsf d_j)+h_ih_j\mathsf P\big],\\
 \mathsf Z_i\big[&\sigma(\mathsf Y_i+\mathsf R_i+
                        \mathsf d_i+\mathsf P)
                         +h_i(\mathsf Z_0+\mathsf Z_1)\big].
\end{aligned}
\end{equation}
Their transport directions are $j$ and $i$, respectively.
The remaining currents, in directions $i$, $1$, $1$, and $i$, have
costs
\begin{equation}\label{dn:high-other-costs}
\begin{aligned}
 &C_M\mathsf z_i h_i(\mathsf z_i+\mathsf R_i+
                                      \mathsf d_i+h_i\mathsf P),\\
 &C_M\mathsf D[\sigma(\mathsf Z_0+\mathsf Z_1+
                                      \sigma\mathsf P)+h_1\mathsf D],\\
 &C_M\mathsf P_\gamma\mathsf R_1+C_Mh_1\mathsf P_\gamma^2,
 \qquad C_M\mathsf P_i(\mathsf Y_i+\mathsf R_i+
                                      \mathsf d_i+h_i\mathsf P).
\end{aligned}
\end{equation}
For the first line of \Cref{dn:high-other-costs}, the high relative
lapse factor multiplies the two low shapes and therefore has
coefficient $h_i^2$.  The shift source is
$q\mathcal A\gamma^{-1}(\xi_1-\xi_0)$, which gives the second line.
The last two lines are the metric and torsion transports.  Their
unweighted sources explain why $\mathsf R_1$ and $\mathsf Y_i$
remain there.  Area and norm differentiation of $D^r$ tensors uses
only $\chi_i=O(h_i)$ on their contracted slots; the scalar weight
uses $L_i\log a=-\kappa+O(h_i)$.  Thus the constant part is favorable
damping and the remaining terms are among the displayed costs.
These placements account for every same-order interaction.

For use on a fixed finite slab, boundedness of $h_i$ already suffices:
substitute \Cref{dn:high-hodge} and multiply all currents by
$e^{-2Kt}$, with $K>C_M$ for the bounded low coefficient set.
Both directions gain damping $K$, giving finite flux bounds.
Thus the estimates are a priori estimates on existing subslabs;
their derivation does not require continuation or exhaustion.

These fixed-slab bounds ensure finiteness, but interpolation later needs
more: high derivatives must cost only an arbitrarily small exponential
rate in $T$. We now obtain that rate by rescaling the one-way limiting
couplings and assigning stronger damping to the bounded optical direction.

For the sharper rate fix $0<\delta\le1$.  The upper-bounded direction
is $d=0$, and the long direction is $l=1$.  When
$\sigma=h_0=h_1=0$, the only nonzero arrows in
\Cref{dn:high-pair-costs,dn:high-other-costs}, after substitution of
\Cref{dn:high-hodge}, are
\[
 \mathsf Y_i,\mathsf z_i\longrightarrow\mathsf P_i,
 \qquad \mathsf D\longrightarrow\mathsf P_0,
 \qquad \mathsf Y_1,\mathsf z_1\longrightarrow\mathsf P_\gamma.
\]
There is also an asymmetric regime.  Give the along-$0$ variables
$\mathsf Y_1,\mathsf Z_0,\mathsf z_0,\mathsf P_0$ arbitrarily
large damping and temporarily remove every product containing one
of them.  At $\sigma=h_1=0$, with $h_0$ bounded, the remaining
arrows are exactly
\[
 \mathsf z_1\longrightarrow\mathsf Y_0,
 \qquad\mathsf z_1\longrightarrow\mathsf P_1,
 \qquad\mathsf z_1\longrightarrow\mathsf P_\gamma.
\]
For example, the $\mathsf P_1$ source contains
$\mathsf Y_1$, which has been removed, rather than $\mathsf Y_0$;
and $\mathsf d_1=0$.  The retained trace and shift equations have
no unsuppressed source returning to $\mathsf z_1$.

Give each paired current multiplier one, the trace currents
multiplier $\Lambda^2$, and the $\mathsf P$ currents multiplier
$\epsilon_P^2$.  Take $\Lambda$ large and $\epsilon_P$ small,
depending only on $M,\delta$ and the bounded low coefficients.
Both displayed limiting quadratic forms then have norm at most
$\delta/4$ in the rescaled energy.  After those choices, continuity
on the bounded low coefficient set gives a fixed threshold
$\varepsilon_*$: the whole form has cost at most $\delta/2$ when
$\sigma,h_0,h_1\le\varepsilon_*$, and the retained long-direction
form has that cost when $\sigma,h_1\le\varepsilon_*$.

Use the common Lipschitz weight
\begin{equation}\label{dn:high-weight}
 \mathfrak w_T=
 \exp\{-2\delta(u+v)-2K_0(u+B_{\mathrm{cut}})_+
                         -2K_0\min(t,S_0)\}.
\end{equation}
Uniform tail smallness from \Cref{dn:uniform-envelopes} permits
$B_{\mathrm{cut}},S_0$ independent of large $T$ such that the
all-small regime holds for $t>S_0,u<-B_{\mathrm{cut}}$ and the
asymmetric regime holds for $t>S_0,-B_{\mathrm{cut}}\le u\le u_2$.
Its directional damping is, almost everywhere,
\[
 -L_k\log\mathfrak w_T
 =2\delta+2K_0\mathbf1_{\{k=0, u>-B_{\mathrm{cut}}\}}
                         +K_0\mathbf1_{\{t<S_0\}}.
\]
If $C_{\mathrm{bd}}$ bounds the complete rescaled form, the products
containing an along-$0$ factor are absorbed by
$K_0>C(C_{\mathrm{bd}}+C_{\mathrm{bd}}^2/\delta)$.
In $t<S_0$ both directions have this damping.  The choice order is
\[
 (\Lambda,\epsilon_P),\quad\varepsilon_*,\quad
                  (B_{\mathrm{cut}},S_0),\quad K_0.
\]
In particular choosing $K_0$ does not change a previously selected
tail threshold.  All constants use only the dependencies stated in
the Lemma.

Integration gives positive null and slice fluxes bounded by the
initial fluxes.  The tilted sides of \Cref{dn:domain} have the same
outgoing signs.  Since $u\le u_2$ and $v\le C(T+1)$,
\[
 1\ge\mathfrak w_T\ge
  \exp\{-C\delta T-2K_0(u_2+B_{\mathrm{cut}})_+-2K_0S_0-C\delta\}.
\]
Removing the weight loses an arbitrarily small exponential rate.
Initial integration over intervals of length $O(T)$ costs only a
polynomial.  In particular the transverse fluxes satisfy
\begin{equation}\label{dn:high-transverse}
 \sup_{y_j}\|\mathsf Y_i+\mathsf d_i\|_{L^2(\dd y_i)}
 +\sup_{y_i}\|\mathsf Z_i\|_{L^2(\dd y_j)}
                         \le e^{o(T)}.
\end{equation}
The intervals here are the available intersections with the initial
partial slab.  Every past coordinate segment stays inside it and
reaches $t=0$, because the side defining functions are nondecreasing
along both $L_0,L_1$.

The weighted angular fluxes are now controlled. Unweighted fields and
ordinary longitudinal jets remain to be recovered. Recover the former
by transport, rather than dividing the
$\mathsf Z$ estimate by $\sigma$.  The damped trace equation and
\Cref{dn:high-hodge} give
\begin{equation}\label{dn:high-trace-convolution}
\begin{split}
 \mathsf z_i(y)\le{}& e^{o(T)}e^{-\kappa(y_i+y_j)}\\
 &+C_M'\int_{-y_j}^{y_i}e^{-\kappa(y_i-r)}h_i
       (\mathsf Y_i+\mathsf d_i+h_i\mathsf P)(r,y_j)\,\dd r.
\end{split}
\end{equation}
The homogeneous $C_Mh_i$ factor has bounded directional integral.
To see explicitly how the trace enters unweighted transport, fix
$y_j$ and integrate \Cref{dn:high-trace-convolution} in its terminal
variable.  Reversing the integration triangle gives
\[
 \int_{-y_j}^{y_i}\mathsf z_i(s,y_j)\,\dd s
 \le e^{o(T)}+\frac{C_M'}{\kappa}
     \int_{-y_j}^{y_i}h_i
       (\mathsf Y_i+\mathsf d_i+h_i\mathsf P)(r,y_j)\,\dd r .
\]
The first transverse flux in \Cref{dn:high-transverse} bounds the
$\mathsf Y_i+\mathsf d_i$ integral by Cauchy--Schwarz, with only a
polynomial interval factor.  Since $h_i$ is uniformly bounded,
$h_i^2\le C h_i$.  The Hodge estimate for $\mathsf R_i$ consequently
gives
\begin{equation}\label{dn:high-pure-integrals}
 \int_{-y_j}^{y_i}(\mathsf z_i+\mathsf R_i)(r,y_j)\,\dd r
 \le e^{o(T)}
       +C_M\int_{-y_j}^{y_i}h_i(r,y_j)\mathsf P(r,y_j)\,\dd r .
\end{equation}
The norm inequalities for metric and torsion transport are
\[
\begin{aligned}
 \mathsf P_\gamma(y)
   &\le e^{o(T)}
       +C_M\int_{-u}^{v}(\mathsf R_1+h_1\mathsf P_\gamma)(u,r)\,\dd r,\\
 \mathsf P_i(y)
   &\le e^{o(T)}
       +C_M\int_{-y_j}^{y_i}
          (\mathsf Y_i+\mathsf R_i+\mathsf d_i+h_i\mathsf P)(r,y_j)\,\dd r .
\end{aligned}
\]
Substitution of \Cref{dn:high-pure-integrals} and the transverse flux
therefore gives
\begin{equation}\label{dn:high-volterra}
 \mathsf P(y)\le e^{o(T)}+
  C_M''\sum_i\int_{-y_j}^{y_i}h_i(r,y_j)\mathsf P(r,y_j)\,\dd r.
\end{equation}
Multiply by $e^{-C\Phi_T(y)}$ and take a supremum.  Along either path,
\[
 \int_{-y_j}^{y_i}h_i(r,y_j)
 e^{-C(\Phi_T(y)-\Phi_T(r,y_j))}\,\dd r\le C^{-1}.
\]
Choose $C>4C_M''$ and absorb the two integrals.  The uniform bound
on $\Phi_T$ proves $\mathsf P,\mathsf z_i\le e^{o(T)}$.
\Cref{dn:high-pure-integrals} now bounds the pure-direction integral
of $\mathsf R_i$ by $e^{o(T)}$ as well.

We record the angular order obtained by the last transport step.
The pointwise $\mathsf P$ bound controls $\gamma$ in $H^{M+1}$ and
$\xi_i$ in $H^M$.  Tame differentiation of angular curvature and
connection therefore gives
\[
 |a\operatorname{Sym}\nabla\xi_j|_{M-1}
       +|aK_\gamma\gamma|_{M-1}+|a\xi_j^{\otimes2}|_{M-1}
 \le C_M q\mathsf P.
\]
Commuting the cross equations with $D$ through order $M-1$ gives
top homogeneous coefficient $C_Mh_j$ on the $\chi_i,\ell_i$ norms.
The other shape and shift terms are bounded by
$C_Mh_i(\mathsf R_j+\mathsf d_j)+C_Mh_ih_j\mathsf P$.
Their integrals in direction $j$ are $e^{o(T)}$ by
\Cref{dn:high-transverse,dn:high-pure-integrals}.  Gronwall thus gives
pointwise $H^{M-1}$ bounds for $\chi_i,\ell_i$.  Finally
$[\partial_v,D^r]b=0$ and the shift source has $H^M$ norm at most
$C_Mq\mathsf P$; its norm derivative adds only $C_Mh_1|b|_M$.
Integration gives a pointwise $H^M$ bound for $b$.
Thus the order-$M$ current gives pointwise $H^{M-1}$ bounds for all
relative fields, and running it with $M\ge A+1$ supplies an angular
target order $A$.  The shift step uses the pointwise torsion bound:
the transverse integration direction of the $\mathsf Z_1$ flux in
\Cref{dn:high-transverse} is $0$.

We next specify the longitudinal induction and its finite angular
buffer.  Fix a largest longitudinal order $n_{\parallel}$ and a
desired final angular order $A_{\rm out}\ge5$, and put
\[
 A_k=A_{\rm out}+2(n_{\parallel}-k),
                    \qquad 0\le k\le n_{\parallel}.
\]
At level $k$ the assertion is an $e^{o(T)}$ bound in reference
$H^{A_k}$ for every ordered transport jet of the relative fields
listed in the Lemma with exactly $k$ longitudinal derivatives: use
$L_i$ on scalars and $\mathfrak L_i$ on angular tensors.  For $b$ we
also keep its ordinary longitudinal
jets at that order.  All earlier levels retain their larger
$H^{A_\ell}$ bounds.  Run the preceding angular argument with
$M\ge A_0+1$ to establish level zero.  The first pure jets in the
bootstrap supply low coefficient bounds; the high-angular induction
begins at $k=1$ from this level-zero conclusion.

For $i\ne j$, the exact identities needed for the induction are
\begin{equation}\label{dn:longitudinal-identities}
\begin{gathered}
 [\mathfrak L_j,\mathfrak L_i]
       =\Lie_{X_{ij}},\qquad X_{ij}=a(\xi_j-\xi_i)^\sharp,\\
 \mathfrak L_i\gamma=2\chi_i,\qquad
 L_i\log\mathcal A=\ell_i+\kappa,\qquad
 \mathfrak L_i(\xi_0+\xi_1)=\dd\ell_i .
\end{gathered}
\end{equation}
For $k\ge1$, Raychaudhuri, torsion transport, and angular curvature
variation give the reductions
\begin{equation}\label{dn:longitudinal-reductions}
\begin{aligned}
 L_i^k\tau_i
   &=L_i^{k-1}(\ell_i\tau_i-|\chi_i|^2),\\
 \mathfrak L_i^k\xi_j
   &=\mathfrak L_i^{k-1}(\beta_i+\tau_i\xi_i),\\
 \mathfrak L_i^k\xi_i
   &=\dd L_i^{k-1}\ell_i-\mathfrak L_i^k\xi_j,\\
 L_i^kK_\gamma
   &=L_i^{k-1}(\divg\beta_i-\tau_iK_\gamma),\\
 [\mathfrak L_j,\mathfrak L_i^k]
   &=\sum_{r=0}^{k-1}
       \mathfrak L_i^r\Lie_{X_{ij}}\mathfrak L_i^{k-1-r}.
\end{aligned}
\end{equation}
The scalar commutator has the same formula with $X_{ij}$ acting by
directional differentiation.  Opposite shape and lapse derivatives
use their cross equations.  The variation of the angular Levi-Civita
connection under $\mathfrak L_i\gamma=2\chi_i$ is
\[
 (\mathfrak L_i\nabla)^C{}_{AB}
 =\gamma^{CD}
   (\nabla_A\chi_{i,BD}+\nabla_B\chi_{i,AD}
                                      -\nabla_D\chi_{i,AB}),
\]
and every derivative of $a$ retains its $a$ factor times relative
fields.  Thus a reduction to a preceding longitudinal level adds at
most two angular derivatives: the largest costs are $\divg\beta_i$
and the outer $\nabla$ on the torsion reduction.  The buffer $A_k$
covers these costs.

Set $U_{i,k}=\mathfrak L_i^k\chi_i$ and $V_{i,k}=L_i^k\ell_i$.
Commuting the cross equations and using
\Cref{dn:longitudinal-reductions} gives
\begin{equation}\label{dn:pure-top-system}
\begin{aligned}
 \mathfrak L_jU_{i,k}
   &=-\tfrac12\tau_jU_{i,k}
       +U_{i,k}\chi_j+\chi_jU_{i,k}+\mathcal F_{i,k},\\
 L_jV_{i,k}
   &=-U_{i,k}\cdot\chi_j+\mathcal G_{i,k}.
\end{aligned}
\end{equation}
Here $\mathcal F_{i,k},\mathcal G_{i,k}$ contain only longitudinal
levels below $k$ and their buffered angular derivatives.  In
particular the first reduction in
\Cref{dn:longitudinal-reductions} replaces the differentiated term
$-\tfrac12(L_i^k\tau_i)\chi_j$ by lower levels.  This is why that
term is absent from the displayed top system.  The new lapse jet
does not occur on its right side.

For completeness, commute the first equation of
\Cref{dn:pure-top-system} with $D$ in increasing angular order.
For $r\ge4$ the tame product and shift commutator inequalities give
\[
\begin{aligned}
 \frac{\dd}{\dd y_j}|U_{i,k}|_r
 &\le C_rh_j|U_{i,k}|_r
   +C_r\mathcal C_{j,r}\sum_{\ell<r}|U_{i,k}|_\ell
   +|\mathcal F_{i,k}|_r,\\
 \mathcal C_{j,r}
 &=1+|\gamma|_{r+1}+|\gamma^{-1}|_{r+1}
                         +|\chi_j|_r+|B_j|_{r+1}.
\end{aligned}
\]
This norm inequality is understood in its regularized form at zero.
At the first four angular orders the low $H^4$ shape and $H^5$ shift
bounds give the same top coefficient directly.  At higher orders
$\mathcal C_{j,r}$ multiplies already controlled lower angular
orders of $U_{i,k}$; it is a source in this subsidiary induction.
It is $e^{o(T)}$ by level zero.  The lower longitudinal induction
bounds $\mathcal F_{i,k}$ using $A_\ell\ge A_k+2(k-\ell)$.
Consequently Gronwall exponentiates only the uniformly bounded
directional integral of $C_rh_j$.  Finite products of earlier
$e^{o(T)}$ bounds and integration over length $O(T)$ remain
$e^{o(T)}$.  After $U_{i,k}$ is controlled, the second equation has
the same transport commutators and the known source
$U_{i,k}\cdot\chi_j$, so it controls $V_{i,k}$ in the same way.

The reductions also cover mixed words.  In an ordered word of length
$k$ acting on a torsion, apply its innermost transport derivative
first: the torsion equation applies to the opposite torsion, and the
torsion-sum identity gives the derivative of the same-index torsion.
The remaining $k-1$ operators then act on $\beta_i,\ell_i$ and
products of fields.  Expanding their angular commutators uses the
same two-derivative buffer.  The metric and relative lapse words
reduce in the same way by \Cref{dn:longitudinal-identities}.
For a mixed shape or lapse word, commute a cross derivative inward;
the commutator in \Cref{dn:longitudinal-reductions} has shorter words,
and the cross equation reduces the resulting innermost derivative.
This gives the already controlled pure or lower levels.  For the shift,
differentiate its exact direction-$1$ equation $k$ times in $u$ and
integrate to obtain $\partial_u^kb$.  Conversion of its source from
transport to ordinary derivatives uses only longitudinal orders
below $k$ of $b$: expanding
$\partial_u=\mathfrak L_0-\Lie_b$ on tensors, or
$\partial_u=L_0-b\cdot\dd$ on scalars, differentiates $b$ at most
$k-1$ times.  Each replacement adds at most one angular derivative,
which is covered by the same buffer.  Mixed shift jets follow
directly from its equation.  Finally
$q^{-1}\partial_vb=\mathcal A\gamma^{-1}(\xi_1-\xi_0)$ is algebraic.
This completes level $k$ and the conversion to ordinary jets.

The required initial pure jets are finitely many entry jets.  They
can also be expressed using tangential $x$ jets and the cross
equations on the noncharacteristic cylinder, where $q=q_0>0$.
For an output order $n$, take $n_{\parallel}=n$ and
$A_{\rm out}=\max(n,5)$, then choose $M\ge A_0+1$ as above.
Include the derivatives defining $K,\beta$ and the initial cross
identities and take their finite maximum as $r(n)$.  All sources in
the argument are polynomial in the finitely many earlier high norms;
only low directional integrals enter an exponential.  This gives
the asserted uniform dependence and leaves $r(n)$ independent of
later frequency-accuracy choices.  The Lemma follows.
\end{proof}

\subsection{The unweighted low error system}

Write $\Delta f=f'-\bar f$.  The following sizes on a sphere use
angular norms in uniformly equivalent low geometry; for the paired
energies we use the primed metric exactly:
\begin{equation}\label{dn:low-sizes}
\begin{aligned}
 Y_i&=|\Delta\beta_i|_0,&
 Z_i&=\sigma|(\Delta Q_i,\Delta m_i)|_0,&
 z_i&=|\Delta w_i|_1,\\
 N_i&=|\Delta\chi_i|_0+|\Delta\ell_i|_0,&
 D_b&=|\Delta b|_1,&R_i&=|\Delta\chi_i|_1,\\
 P&=|\Delta\gamma|_1+|\Delta\log\mathcal A|_0
                              +|\Delta\xi_0|_0+|\Delta\xi_1|_0.
\end{aligned}
\end{equation}
A squared sum below means the sum of the component squares, with
each component assigned its indicated transport direction.  This
convention matters for the components of $P$.
Set
\[
 V=P+D_b+\sum_i(N_i+Y_i+Z_i+z_i),\qquad
                             H_s=\sigma+h_0+h_1.
\]

\begin{proposition}[Quadratic low error]\label{prop:dn-low-error}
For every $\eta_2>0$, the input loss $\eta$ and the support fraction
$\nu$ may be chosen sufficiently small, with the finite base orders
fixed, so that on every existing bootstrap subslab
\begin{equation}\label{dn:low-conclusion}
 \sup_{0\le t\le L_T}
 \| (\Delta\gamma,\Delta\log\mathcal A,\Delta b,
                   \Delta\chi_i,\Delta\ell_i,\Delta\xi_i)
                         \|_{L^2(I_T(t)\times\Sph)}
 \le \amp^2e^{\eta_2T}.
\end{equation}
The supremum is restricted to available slices. At the same times,
\begin{equation}\label{dn:low-auxiliary-conclusion}
 \sum_i\|Y_i+Z_i+z_i\|_{L^2(I_T(t))}
 +\|(\nabla\Delta\gamma,\nabla\Delta b)
                  \|_{L^2(I_T(t)\times\Sph)}
 \le \amp^2e^{\eta_2T}.
\end{equation}
\end{proposition}

\begin{proof}
We give the operator differences first, to check derivative counts
before estimating the currents.  Put $h_{AB}=\Delta\gamma_{AB}$
in the next two displays only.  The difference of angular connections
is exactly
\begin{equation}\label{dn:connection-difference}
 C^C{}_{AB}:=(\Gamma')^C{}_{AB}-(\bar\Gamma)^C{}_{AB}
 =\tfrac12(\gamma'^{-1})^{CD}
 (\bar\nabla_Ah_{BD}+\bar\nabla_Bh_{AD}-\bar\nabla_Dh_{AB}).
\end{equation}
Thus $|C|_0\le C|\Delta\gamma|_1$.  For a barred one-form $f$,
\begin{equation}\label{dn:operator-differences}
\begin{aligned}
 (\divg'-\overline{\divg})f
 &=((\gamma'^{-1})^{AB}-(\bar\gamma^{-1})^{AB})\bar\nabla_Af_B
            -(\gamma'^{-1})^{AB}C^D{}_{AB}f_D,\\
 (\curl'-\overline{\curl})f
 &=\tfrac12((\epsilon')^{AB}-\bar\epsilon^{AB})(\dd f)_{AB}.
\end{aligned}
\end{equation}
The analogous divergence formula on two-tensors has one connection
difference for each covariant slot.  No second derivative of $h$
occurs.  Against $\bar\beta_i$ or $\bar\chi_i$, these differences
cost $Ch_iP$, since their required low angular derivatives are
$O(h_i)$; against bounded barred torsion they cost $CP$.

The exact torsion-sum identity gives
\begin{equation}\label{dn:H-one}
 |\Delta\log\mathcal A|_1\le CP,
                   \qquad |\Delta\mathcal A|_1\le CP.
\end{equation}
For the trace conversion, use the exact identity
\begin{equation}\label{dn:trace-conversion}
 \Delta\tau_i=\mathcal A'\Delta w_i+
                      (\Delta\mathcal A)\bar w_i,
 \qquad |\Delta\tau_i|_1\le C(z_i+h_iP).
\end{equation}
The trace of $\Delta\chi_i$ in the primed metric is
$\Delta\tau_i-(\gamma'^{-1}-\bar\gamma^{-1})\bar\chi_i$.
Writing $C_{\rm inv}=\gamma'^{-1}-\bar\gamma^{-1}$, the full
shape Hodge datum is exactly
\[
 \divg'\Delta\chi_i-\dd\tr'\Delta\chi_i
 =\Delta\beta_i-(\divg'-\overline{\divg})\bar\chi_i
                              +\dd(C_{\rm inv}:\bar\chi_i).
\]
Apply the primed trace-free divergence estimate, including the
undifferentiated tensor norm to cover its kernel, and then the
primed one-form div--curl estimate.  \Cref{dn:operator-differences}
and \Cref{dn:trace-conversion} yield
\begin{equation}\label{dn:difference-hodge}
 R_i\le C(Y_i+N_i+z_i+h_iP),\qquad
 |\Delta\xi_j|_1\le C(\sigma^{-1}Z_i+P).
\end{equation}
For the latter, the divergence data are
$\Delta K_\gamma-\Delta m_i$ and the curl datum is $\Delta c_i$,
plus the displayed coefficient differences on $\bar\xi_j$.
Precisely,
\[
 \divg'\Delta\xi_j=\Delta K_\gamma-\Delta m_i
                  -(\divg'-\overline{\divg})\bar\xi_j,
 \qquad
 \curl'\Delta\xi_j=\Delta c_i
                  -(\curl'-\overline{\curl})\bar\xi_j.
\]
These first-order elliptic constants need only the uniform low
Lipschitz geometry.  The estimate does not assert an unweighted
bound for the torsion derivative by $Z_i$.

When subtracting a transport equation, keep the actual operator:
\begin{equation}\label{dn:transport-differences}
\begin{aligned}
 \mathfrak L_i'f'-\bar{\mathfrak L}_i\bar f
    &=\mathfrak L_i'\Delta f+\Lie_{\Delta B_i}\bar f,\\
 L_i'f'-\bar L_i\bar f&=L_i'\Delta f+
                                          \Delta B_i\cdot\dd\bar f.
\end{aligned}
\end{equation}
For a one-form, the last Lie term consists of
$\Delta B_i\cdot\partial\bar f+(\partial\Delta B_i)*\bar f$.
It costs $Ch_iD_b$ on $\bar\beta_i,\bar\chi_i$, $CD_b$ on
bounded barred torsion, and $CD_b$ on $\bar Q_i,\bar m_i$ before
their $\sigma$ weight.  On $\bar\ell_i=-\kappa+O(h_i)$ it costs
$Ch_iD_b$, since the constant $-\kappa$ has zero angular derivative.

\paragraph{The paired currents.}
Writing $k=\Delta K_\gamma$ and $c=\Delta c_i$, the primed
principal part of the difference system is
\[
 \mathfrak L_j'\Delta\beta_i
        =\tfrac{a'}2(\dd k+J'\dd c)+F_i,
 \qquad
 L_i'\Delta Q_i=(\divg',\curl')\Delta\beta_i+G_i,
 \qquad L_i'\Delta m_i=G_i^m.
\]
There are no angular derivatives in the last principal equation.
Use the currents with densities
\begin{equation}\label{dn:pair-densities}
 e_{Y_i}=\tfrac12|\Delta\beta_i|_{\gamma'}^2
                 \quad\hbox{in direction }j,
 \qquad
 e_{Z_i}=\tfrac{a'}4
       (|\Delta Q_i|^2+|\Delta m_i|^2)
                 \quad\hbox{in direction }i.
\end{equation}
The ratio of scalar to one-form weight is exactly $a'/2$.
Since $\dd$ and $J'\dd$ have adjoints $-\divg'$ and $-\curl'$,
their principal contribution, integrated on the full sphere, is
\begin{equation}\label{dn:principal-cancellation}
 -\tfrac12\int_{\Sph}
 \langle\Delta\beta_i,
       (\dd a')\,k+J'(\dd a')\,c\rangle_{\gamma'}\,\dd\mu_{\gamma'}.
\end{equation}
This has size $C\sigma Y_iZ_i$, because
$\dd a'=a'(\xi_0'+\xi_1')$ and $a'\asymp q$.
The derivative of the scalar weight contains
$L_i'a'=a'\ell_i'=a'(-\kappa+O(h_i))$, whose constant part is
favorable damping.  The area and tensor-norm derivatives contain
only low shapes and shifts.  They cost
$C(h_jY_i^2+h_iZ_i^2)$.
Indeed, after the angular transport divergence is integrated out,
their exact contributions to the two densities are
\[
 \tfrac{\tau_j'}2|\Delta\beta_i|^2
       -\chi_j'((\Delta\beta_i)^\sharp,(\Delta\beta_i)^\sharp),
 \qquad
 \tfrac{a'}4(\ell_i'+\tau_i')
                         (|\Delta Q_i|^2+|\Delta m_i|^2).
\]

Here are all nonprincipal placements in $F_i,G_i,G_i^m$.
The change of the angular principal coefficient on the barred field
is
\begin{equation}\label{dn:principal-coefficient-error}
 \tfrac12\{(a'-\bar a)(\dd\bar K+\bar J\dd\bar c_i)
                      +a'(J'-\bar J)\dd\bar c_i\}.
\end{equation}
It costs $CqP$, since $\Delta a=q\Delta\mathcal A$ and $J'-\bar J$
is an algebraic metric difference.  The divergence and curl
coefficient differences against $\bar\beta_i$ cost $Ch_iP$ and
therefore $C\sigma h_iP$ after scalar weighting.  Shift differences
cost $Ch_iD_b$ in the first equation and $C\sigma D_b$ in the other
two.  The remaining $a$ terms cost at most
$C\sigma(Z_0+Z_1+P)$ in the first equation.  Indeed they consist of
$a'\Delta K$, $a'\nabla'\Delta\xi$, bounded low factors times
$a'\Delta\xi$, and coefficient differences on bounded barred
factors.  The shape products have bound
\begin{equation}\label{dn:shape-product-error}
 C(h_iR_j+h_jR_i+h_ih_jP).
\end{equation}
In particular the angular connection difference in
$\chi_i*\nabla\chi_j$ has the explicit factor
$\chi_i'*C*\bar\chi_j$ and costs $Ch_ih_jP$.
In the scalar equations, the weighted remainder has bound
\[
 C\{\sigma(R_i+D_b+P)+h_i(Z_0+Z_1)\}.
\]
These lists exhaust the factors of \Cref{dn:D3}. Put $Z=Z_0+Z_1$.
In particular all
uses of the apparently dangerous second estimate in
\Cref{dn:difference-hodge} occur with one of the factors
\begin{equation}\label{dn:compensation}
 a'|\Delta\xi|_1\le C(\sigma Z+qP),\qquad
 \sigma h_i|\Delta\xi|_1\le C(h_iZ+\sigma h_iP).
\end{equation}
There is no naked $\sigma^{-1}Z$ in a current source.
The combined remaining paired cost is consequently bounded by
\begin{equation}\label{dn:low-pair-cost}
 C H_s(Y_i+Z_i)(V+R_0+R_1),
\end{equation}
apart from the defects paired with the current variables.

\paragraph{Trace, shape, lapse, and shift currents.}
Subtract \Cref{dn:damped-trace} with the actual $L_i'$:
\begin{equation}\label{dn:trace-difference-equation}
 (L_i'+\kappa)\Delta w_i
 =-\Delta(\mathcal A^{-1}|\chi_i|^2)
          -\Delta B_i\cdot\dd\bar w_i+\mathcal F_i^w.
\end{equation}
The defect $\mathcal F_i^w$ has the input bound.  Its one angular
commutation has cost at most
$Ch_i(z_i+R_i+D_b+h_iP)$, in addition to the homogeneous damping:
the differentiated quadratic shape has a low $h_i$ factor,
$\Delta\mathcal A$ is controlled in $H^1$ by \Cref{dn:H-one},
and the shift term uses $\bar w_i$ through two low angular
derivatives, all $O(h_i)$.  The commutator with $L_i'$ uses only
the low first angular derivative of $B_i'$.  Thus the remaining
$z_i^2$ current cost is $CH_sz_i(V+R_i)$; the constant
$2\kappa z_i^2$ has favorable sign.  In particular no
$\Delta\ell_i$ term is generated by this normalized trace equation.

For $N_i$ use the shape and lapse cross equations in direction $j$.
Their source factors are precisely
\[
 a\nabla\xi,\quad aK_\gamma,\quad a\xi^2,
                                      \quad \chi_i*\chi_j.
\]
Their differences are bounded by
$C\{\sigma(Z_0+Z_1)+qP+h_iN_j+h_jN_i+h_ih_jP+h_iD_b\}$.
For example, using the connection difference in
\Cref{dn:connection-difference},
\begin{align*}
 \Delta(a\operatorname{Sym}\nabla\xi_j)
 &=a'\operatorname{Sym}\nabla'\Delta\xi_j
   +(\Delta a)\operatorname{Sym}\bar\nabla\bar\xi_j
   -a'\operatorname{Sym}(C*\bar\xi_j),\\
 \Delta(aK_\gamma\gamma)
 &=a'(\Delta K_\gamma)\gamma'
      +a'\bar K_\gamma\Delta\gamma
      +(\Delta a)\bar K_\gamma\bar\gamma.
\end{align*}
This uses \Cref{dn:compensation} and
\Cref{dn:transport-differences}; the constant part of $\ell_i$
does not contribute.  The squared $N_i$ currents therefore cost
at most $CH_sN_i(V+R_0+R_1)$.  In particular there is no
unsuppressed return from a lapse or metric error to $N_i$.

For the $H^1$ shift current in direction $1$, subtract its exact
equation:
\[
 \partial_v\Delta b=q\Delta\{
             \mathcal A\gamma^{-1}(\xi_1-\xi_0)\}.
\]
One angular derivative uses \Cref{dn:H-one} and
\Cref{dn:compensation}, giving a source
$C(\sigma(Z_0+Z_1)+qP)$ and a squared-current cost
$CH_sD_bV$.  No differentiated shift is needed on its right side.

\paragraph{The unweighted recovery currents.}
Use direction $1$ for $\Delta\gamma$ in $H^1$ and
$\Delta H$ in $L^2$, where $H=\log\mathcal A$:
\begin{equation}\label{dn:P-transport}
 \partial_v\Delta\gamma=2\Delta\chi_1,
                    \qquad\partial_v\Delta H=\Delta\ell_1.
\end{equation}
For comparison, an arbitrary-direction equation would be
$L_i'\Delta H=\Delta\ell_i-\Delta B_i\cdot\dd\bar H$;
the chosen direction has $\Delta B_1=0$.
For $\Delta\xi_j$ use direction $i$ and
\begin{equation}\label{dn:torsion-difference}
 \mathfrak L_i'\Delta\xi_j
 =\Delta\beta_i+\tau_i'\Delta\xi_i
       +(\Delta\tau_i)\bar\xi_i
       -\Lie_{\Delta B_i}\bar\xi_j+\mathcal F_i^\xi.
\end{equation}
In its undifferentiated norm the trace difference costs
$C(N_i+h_iP)$ directly from the definition of $\tau_i$.
The last Lie derivative involves bounded barred torsion and its
first angular derivative, giving $CD_b$.  Use fixed-background
angular norms for these recovery currents to avoid any unnecessary
time differentiation of their norm coefficients.  Their total cost,
including harmless low transport commutators, is
\begin{equation}\label{dn:low-P-cost}
 CP\big\{\sum_i(Y_i+N_i+R_i)+D_b+H_sP\big\}.
\end{equation}
All occurrences of unweighted first angular torsion derivatives
have already been confined to \Cref{dn:compensation}.

\paragraph{The growth rate and the late support.}
Substitute \Cref{dn:difference-hodge} in the preceding costs.
At $H_s=0$, all currents except the $P$ currents have zero
unfavorable cost.  The surviving directed interactions are
\begin{equation}\label{dn:limiting-graph}
 (Y_i,N_i,z_i,D_b)\longrightarrow P.
\end{equation}
There is no reverse arrow.  Give the component $P$ currents a
small fixed multiplier $\delta_P^2$ and every other current fixed
positive multipliers, equal on each Bianchi pair.  Let
$\mathcal E_T(t)$ be the resulting spatial slice energy.  Its
density is uniformly equivalent, with constants depending on
$\delta_P$ but not $T$, to
\begin{equation}\label{dn:energy-equivalence}
 \sum_i(Y_i^2+Z_i^2+z_i^2+N_i^2)+D_b^2+\delta_P^2P^2.
\end{equation}
In particular it controls all the unweighted quantities claimed in
\Cref{dn:low-conclusion}.  The factors $1/2$ from the time
components of $L_i'$ merely change this equivalence by constants.

Given a small rate $\rho_*>0$, first take $\delta_P$ small so
that the form in \Cref{dn:limiting-graph}, expressed in scaled
variables, costs at most $\rho_*\mathcal E_T/2$.
Next choose a fixed $h_*>0$ so that every term with factor $H_s$
costs at most $\rho_*\mathcal E_T/2$ when $H_s\le h_*$.
On the remaining bounded low coefficient set the rate is bounded
by a constant $C_*$ independent of $T$.  This constant is fixed
before $\nu$ is chosen.

All difference fields vanish in $v<(1-\nu)T$ by protected
agreement, causal monotonicity and uniqueness of the optical
transports.  Fix a large $R$ so that the background envelope
suprema are small outside $[-R,R]$, and a fixed $t_0$ so that both
$\sigma$ and $q=\sigma^2$ are below their required thresholds for
$t>t_0$.  For large $T$, $h_1$ is small on the support for $t>t_0$.
At those times, the only possible
failure of $H_s\le h_*$ is where $u\ge-R$.  At any such supported
point,
\[
 2t=u+v\ge (1-\nu)T-R,
 \qquad t\le L_T=T/2+O(1).
\]
Consequently the slice times containing a bad supported point lie in
\[
 [0,t_0]\ \cup\
 \left[\frac{(1-\nu)T-R}{2},L_T\right]
\]
after omitting any empty interval.  Their total length is at most
\begin{equation}\label{dn:bad-strip}
                 t_0+\nu T/2+O(1).
\end{equation}
The added $e_T$ in the majorants is uniformly below the threshold
for large $T$.  Thus the integrated homogeneous rate is bounded by
$C\rho_*T+C_*\nu T/2+O(1)$, and is an arbitrarily small multiple
of $T$ after choosing $\rho_*$ and then $\nu$.

Both optical directions have nonnegative outgoing flux on the
tilted sides: the derivatives of $u+c_0t/T$ and $v+c_0t/T$ along
$L_0',L_1'$ are nonnegative, and are strictly positive where
needed.  There are no angular boundary terms because the sphere
is complete.  The defects have magnitude $F_T$ on a spatial
interval of length $O(T)$ and therefore add at most
$C\operatorname{poly}(T)F_T\sqrt{\mathcal E_T}$ to the squared
energy inequality.  The initial discrepancy has the same bound
after the finite spatial integration.  The square-root inequality
(or its regularization at zero) and Gronwall yield
\begin{equation}\label{dn:quadratic-energy}
 \sqrt{\mathcal E_T(t)}\le
 C\operatorname{poly}(T)\amp^2e^{\eta T}
 \exp\{C\rho_*T+C_*\nu T/2+O(1)\}
                     \le\amp^2e^{\eta_2T}.
\end{equation}
Choose the several positive losses strictly below $\eta_2$ to
absorb the fixed and polynomial factors.  This proves the Proposition.
\end{proof}

\subsection{Interpolation, improvement, and continuation}

\begin{proof}[Proof of \Cref{thm:dn-comparison}]
We first make explicit the finite order choices.  Choose a spatial
target order $k\ge2s$, enlarged if necessary to improve every
bootstrap and continuation norm.  The time recovery below uses at
most two spatial derivatives for each time derivative.  Choose an
integer $M$ much larger
than $k+2$ so that, with $\vartheta=(k+2)/M$,
\begin{equation}\label{dn:interpolation-choice}
             2\alpha(1-\vartheta)>(1+c_1)\alpha
                         \quad\text{for some }0<c<c_1<1/2.
\end{equation}
The constant $2$ in $k+2$ supplies Sobolev embedding on the
three-dimensional slices.  Here $M$ is the requested slice Sobolev
order; the proof of \Cref{lem:dn-high-uniform} runs its angular
currents at the larger order specified there.  That Lemma determines
a finite initial buffer $r(M)$ and its finite template/exterior
jet requirements.  Include the finite order $m$ in that base
budget.  Choose the rate tolerance $\eta_2$ small enough that
\Cref{dn:interpolation-choice} retains a strict gap after adding
all interpolation and later time-recovery losses.  In
\Cref{prop:dn-low-error} choose the fixed current multipliers,
the small-region threshold, the growth tolerance, then $\nu$, and
finally sufficiently small input/exterior losses.

Let $P_{\mathrm{base}}\ge1$ be the maximum of the finitely many
positive frequency exponents required at these fixed base orders.
Choose $\zeta>0$ so small that each has its allocated loss, in
particular $\zeta P_{\mathrm{base}}<\alpha/2$. Only now increase the
finite hierarchy and beam accuracy orders to make all linear
defects as small as required, including the finite inverse-lapse
losses in deriving linearized double-null identities.  The
base-loss statement in \Cref{prop:packet-preparation} ensures
that this last operation changes constants and required fixed
background jets, but does not increase any previously budgeted
positive power bounded by $P_{\mathrm{base}}$ of $\lambda$.
Finally take $T$ large.
Every high input is a finite entry norm already supplied by the
preparation.  This order of choices has no unknown-high-norm
smallness requirement.

Local smooth vacuum theory and noncharacteristic eikonal and
angle transport start the primed solution at the actual cylinder
data.  On each available slice, \Cref{dn:low-conclusion} and
\Cref{lem:dn-high-uniform} give, for each component error $f$,
\[
 \|f\|_{L^2}\le e^{-(2\alpha-\eta_2)T},
                       \qquad\|f\|_{H^M}\le e^{o(T)}.
\]
For the second inequality the barred fields have their individual
high bounds directly from the linear construction.  Sphere charts
and extension in the interval variable give
\begin{equation}\label{dn:spatial-interpolation}
 \|f\|_{C^k(I_T(t)\times\Sph)}
 \le C\operatorname{poly}(T)
      \|f\|_{L^2}^{1-\vartheta}\|f\|_{H^M}^{\vartheta}
                    \le e^{-(1+c_1)\alpha T}
\end{equation}
after reducing $c_1$ within its fixed margin.  The lower bound for
the interval length ensures uniform local extension constants;
covering a longer interval has at worst a polynomial cost.
This argument is performed on each slice and needs no positive
lower bound for the available time thickness.

To obtain time derivatives, use $\partial_v=(\partial_t+
\partial_x)/2$ in \Cref{dn:recovery-equations}:
\[
\begin{aligned}
 \partial_t\gamma&=4\chi_1-\partial_x\gamma,\\
 \partial_t\log\mathcal A
   &=2(\ell_1+\kappa)-\partial_x\log\mathcal A,\\
 \partial_tb
   &=2q\mathcal A\gamma^{-1}(\xi_1-\xi_0)-\partial_xb .
\end{aligned}
\]
Use the shape/lapse cross equations and the torsion equations for the
other fields.  For a tensor
transported in direction $0$ use
$\mathfrak L_0=(\partial_t-\partial_x)/2+\Lie_b$; for a scalar
replace the Lie derivative by $b\cdot\dd$.  Thus each time
derivative is twice the corresponding right side, plus a spatial
derivative and, when appropriate, an angular Lie term.  A shape
or lapse cross equation contains at most two angular derivatives
of the metric and one of the torsion; the torsion equation contains
one angular shape derivative.  Thus $r$ time derivatives and $p$
spatial derivatives require at most spatial order $p+2r$.  The choice
$k\ge2s$ covers every requested ordinary derivative of total order
at most $s$.  Subtract the
equations with the actual principal operators, as above.  All
products use already controlled lower error jets and individual
jets from \Cref{lem:dn-high-uniform}; barred defects have the
prescribed finite differentiated bounds.  There is no division by
$q$.  Finite product losses can be assigned arbitrarily small
rates in that Lemma, leaving \Cref{dn:main-error} with $c<c_1$.
The relative shift conclusion follows by differentiating the
algebraic identity
\[
 q^{-1}\partial_vb=\mathcal A\gamma^{-1}(\xi_1-\xi_0).
\]

The template change at the fixed bootstrap orders is at most
$e^{-(\alpha-\eta_0)T}$ with $\eta_0<\alpha/2$.
Together with \Cref{dn:main-error}, this is strictly smaller
than $\tfrac12 e^{-\alpha T/2}$ for large $T$.  It improves
\Cref{dn:bootstrap}, preserves the nondegeneracy margins, and
provides the finite Sobolev buffers required for restart.

We finish the existence argument, since an a priori comparison on
a putative full domain would not suffice.  On each fixed finite
experiment, $q\ge q_0e^{-2\kappa L_T}>0$.  All requisite finite
norms and inverse metrics are bounded up to a possible terminal
slice.  Use the finite-domain restart in
\cite[\CompanionRef{foundation}{in:continuation}, proof]{KerrCTwo} as follows.
On an earlier slice take its induced metric and second fundamental
form, initialize wave gauge, and extend the smooth slice jets
across the interval endpoints on an enlarged margin.  Their
required finite Sobolev norms and positivity have controlled
extensions.  Constraints are needed only on the original margin.
The reduced wave equation has a positive local existence time
depending on those finite norms and the finite positive lapse.
Solve the eikonal pair and angular transport from the slice with
their inherited values.  Transversality and invertibility persist
on a smaller margin for a time bounded below at this fixed $T$.

Restrict the new solution to the region with exiting sides.
There $\dd u,\dd v\ge0$ on future causal vectors, so artificial
extension data cannot enter the retained domain.  Constraint and
gauge propagation give the vacuum equations there.  Local
geometric uniqueness, followed by uniqueness for the eikonal and
angle transports, identifies this solution with the existing one
on their overlap.  Restart before any proposed terminal time;
the uniform bounds at that fixed $T$ contradict termination.
Slightly enlarged intervals and sphere patches give smooth
endpoint jets throughout.  The restart time need not be uniform
as $T\to\infty$.

Finally, smoothness at every order is preserved.  At each fixed
$T$ the fixed-region version of the high-current proof, with the
weight $e^{-2Kt}$, propagates any larger finite order from the
smooth entry data.  Its constants may depend on that order and
on $T$; no simultaneous subexponential estimate at all orders
is asserted or needed.  The constructed solution therefore
reaches all of \Cref{dn:domain} smoothly.  On the initial collar
it is the same vacuum solution as the prepared exterior solution
in its newly initialized optical coordinates, by the same local
uniqueness.  The whole-bridge attachment is carried out in
\Cref{prop:experiment-attachment}.  This proves the Theorem.
\end{proof}
